\documentclass[a4paper, twoside, openright]{article}
\usepackage[T1]{fontenc} % Font encoding, T1 = it
\usepackage[utf8]{inputenc} % Input encoding - per caratteri particolari
\usepackage[english,italian]{babel} % Lingua principale italiano, con parti in inglese

\usepackage{graphicx} % Per includere immagini esterne
\usepackage{lipsum} % genera testo fittizio
\usepackage[a4paper,top=3cm,bottom=3cm,left=3cm,right=3cm]{geometry} %impaginazione e margini documento
\usepackage[fontsize=12pt]{scrextend} %dimensione font
\raggedbottom % Se la pagina non è completa, lascia lo spazio alla fine
\linespread{1.25}
\usepackage{fancyhdr}

% Package for matrix
\usepackage{amsmath}

\pagestyle{fancy}
\fancypagestyle{plain}{%
  \renewcommand{\headrulewidth}{0pt}%
  \fancyhf{}%
}

% IMMAGINI %
\usepackage{subfig}
\usepackage{graphicx}
\usepackage{float}
\graphicspath{ {./images/} }

% STILE LINK ESTERNI %
\usepackage{xcolor}
\definecolor{linkColor}{RGB}{2,11,120}
\definecolor{printLinkColor}{RGB}{0,0,0}

\usepackage[colorlinks=true, allcolors=printLinkColor]{hyperref}
\newcommand\anchor[2]{%
  \href{#2}{#1}\footnote{\url{#2}}%
}

% FONT
\usepackage{tgschola}

% STILE BLOCCHI DI CODICE %
\definecolor{bgTitleRed}{RGB}{85,45,50}
\usepackage[T1]{fontenc}
\usepackage[ttdefault=true]{AnonymousPro}
\usepackage{listings}
\usepackage{minted}
\usepackage{tcolorbox}
\tcbuselibrary{listings, breakable, minted, skins}
\tcbset{listing engine=minted}

\newtcblisting{bashCode}[2][]{
    breakable,
    listing only, #1 ,title=#2,
    minted language=bash,
    minted style=vs,
    coltitle=white,
    colbacktitle=bgTitleRed,
    toptitle=3mm, bottomtitle=2.5mm,
    top=2mm, bottom=3mm,
    fonttitle=\ttfamily,
    enhanced, frame hidden, 
    minted options={fontfamily=AnonymousPro, 
    tabsize=4, breaklines, autogobble, linenos=false}}

\begin{document}

% Title page
\begin{titlepage}
    \centering
    
    % Logo Università
    \includegraphics[width=7cm]{images/univr.png}\par\vspace{1cm}
    
    % Intestazione
    {\LARGE UNIVERSITÀ DI VERONA \par}
    \vspace{2mm}
    {\large DIPARTIMENTO DI INFORMATICA \par}
    \vspace{5mm}
    {\LARGE Corso di Laurea Triennale in Informatica \par}
    
    \vspace{15mm}
    
    % Titolo tesi
    {\Huge \textbf{Uso di tecniche di LLM per analizzare traiettorie turistiche e per suggerire la prossima attrazione turistica da visitare}\par}
    
    \vfill
    
    % Relatore / Candidato
    \begin{minipage}[t]{0.47\textwidth}
        \raggedright
        {\large Relatrice:}\\[3mm]
        {\Large \textbf{Sara Migliorini}}
    \end{minipage}
    \hfill
    \begin{minipage}[t]{0.47\textwidth}
        \raggedleft
        {\large Candidato:}\\[3mm]
        {\Large \textbf{Mattioli Simone}}
    \end{minipage}
    
    \vspace{20mm}
    
    % Anno accademico
    {\large Anno Accademico 2024--2025}
\end{titlepage}

% Index page
\newpage
\fancyhf{}  % reset header & footer
\fancyhead[RO,LE]{ Indice } %RO=right odd, LE=left even
\tableofcontents

% Footer setting
\newpage
\fancyfoot[RO,LE]{\thepage}
%\fancyfoot[LO,RE]{\thepage}
\setcounter{page}{1}

% Chapters
\section{Introduzione}
\fancyhead{}    % reset header
\fancyhead[RO,LE]{Introduzione}
% PROVA:
%\fancyhead[RO]{Introduzione (1}
%\fancyhead[LE]{1) Introduzione}
La previsione della mobilità umana si è affermata come un'area di ricerca cruciale, con applicazioni che spaziano dalla pianificazione urbana ai sistemi di trasporto, dai sistemi di raccomandazione alla salute pubblica. Gli approcci tradizionali alla previsione della mobilità si sono basati in larga misura su modelli statistici, algoritmi di apprendimento automatico e architetture di deep learning specificamente progettate per dati sequenziali. Tuttavia, i recenti progressi nei Large Language Model (LLM) hanno aperto nuove strade per comprendere e prevedere i modelli di comportamento umano attraverso capacità di elaborazione del linguaggio naturale.

La sfida fondamentale nella previsione della mobilità umana risiede nel catturare la complessa interazione tra fattori spaziali, temporali e comportamentali che influenzano le decisioni di movimento individuali. Gli approcci tradizionali di apprendimento automatico richiedono spesso un'ampia progettazione di feature e conoscenze specifiche del dominio per codificare efficacemente queste relazioni. Al contrario, i LLM possiedono capacità intrinseche di comprensione delle relazioni e dei modelli contestuali, che possono essere sfruttate per prevedere i modelli di mobilità attraverso prompt e informazioni contestuali attentamente progettati.

Questo articolo presenta un framework che sfrutta la potenza dei Large Language Model per la previsione della mobilità umana in contesti turistici ( nello specifico della città di Verona ). Il nostro approccio si differenzia dai metodi convenzionali perché tratta la previsione della mobilità come un compito di comprensione del linguaggio naturale, in cui i modelli di visita storici, le informazioni geografiche e il contesto temporale vengono codificati come prompt testuali per l'inferenza LLM.

I principali contributi di questo lavoro includono:

\begin{itemize}
\item Un nuovo framework per la previsione della mobilità umana utilizzando modelli di linguaggio di grandi dimensioni (LLM) che richiede un'ingegneria minima delle funzionalità
\item Valutazione completa di diversi tipi di informazioni contestuali sull'accuratezza della previsione
\item Implementazione utilizzando modelli open source Llama 3.1, eliminando le dipendenze da API proprietarie
\item Esperimenti approfonditi su dati turistici reali tratti dal dataset VeronaCard
\item Analisi dell'impatto della prossimità geografica e dei modelli temporali sulle prestazioni di previsione della mobilità
\end{itemize}
\newpage
\section{Platooning di veicoli}
\fancyhead{}
\fancyhead[RO,LE]{Platooning di veicoli}

\subsection{Equazioni nel codice}

\subsection{Modello matematico del Platooning}

\subsection{Equazioni in forma matriciale}

\subsection{Campionamento e discretizzazione delle equazioni}
\newpage
\section{Simulazione del Platooning}
\fancyhead{}
\fancyhead[RO,LE]{Simulazione del Platooning}

\subsection{Implementazine della simulazione}

\subsection{Rappresentazione software delle equazioni del Platooning}
\newpage
\section{Frontend}
\fancyhead{}
\fancyhead[RO,LE]{Frontend}

\subsection{Impostazioni}

\subsection{Grafici}

\subsection{Lingue}
\newpage
\section{Esperimenti}
\fancyhead{}
\fancyhead[RO,LE]{Esperimenti}

\subsection{Modello stabile}

\subsection{Modello instabile}
\newpage
\section{Conclusioni}
\fancyhead{}
\fancyhead[RO,LE]{Conclusioni}

\newpage
\input{tex/7-bibliography}
\newpage

% Figures
\fancyhf{}
\fancyhead[RO,LE]{ Elenco delle figure }
\listoffigures
\listoftables
\newpage

% Bibliography
\fancyhf{}
\fancyhead[RO,LE]{ Bibliografia }
\input{tex/7-bibliography}
%\bibliography{tex/bibliography}

% Ringraziamenti
\newpage
\fancyhf{}
\fancyhead[RO,LE]{ Ringraziamenti }
\input{tex/8-thanks-fede}


% End of document
\end{document}


EXAMPLES:

LINK EXAMPLE
\anchor{Name External Link}{https://www.google.com/}

CODE EXAMPLE
\begin{bashCode}{Example Code Block}
    int a = 5;
\end{bashCode}
to insert: %\input{path.tex}

MONOSPACE INLINE TEXT EXAMPLE
\texttt{Example Text}

MONOSPACE TEXT EXAMPLE
\begin{verbatim}
    monospaced text
\end{verbatim}

MATRIX EXAMPLE
$\begin{bmatrix}
    a & b & c\\
    d & f & g
\end{bmatrix}$ · 
$\begin{bmatrix}
    a & b & c\\
    d & f & g
\end{bmatrix}$ = 
$\begin{bmatrix}
    a & b & c\\
    d & f & g
\end{bmatrix}$

IN-LINE EQUATION EXAMPLE
\(x^n + y^n = z^n\)

STAND-ALONE EQUATION EXAMPLE
\[ x^n + y^n = z^n \]

IMAGE EXAMPLE
\begin{figure}[H]
    \centering
    \includegraphics[width=\textwidth]{image.png}
    \caption{caption}
    \label{fig:reference}
\end{figure}

SIDE IMAGES EXAMPLE
\begin{figure}[H]
    \centering
    \begin{subfigure}[b]{0.49\textwidth}
        \centering
        \includegraphics[width=\textwidth]{image1.png}
        \caption{caption1}
        \label{fig:reference1}
    \end{subfigure}
    \hfill
    \begin{subfigure}[b]{0.49\textwidth}
        \centering
        \includegraphics[width=\textwidth]{image2.png}
        \caption{caption2}
        \label{fig:reference2}
    \end{subfigure}
    \caption{total-caption}
    \label{fig:total-reference}
\end{figure}

TABLE EXAMPLE
\begin{table}[!htbp]
    \centering
    \captionsetup{justification=centering}
    \begin{tabular}{|p{4.5cm}|p{9.5cm}|}
        \hline
        a1 & a2\\
        \hline
        b1 & b2 \\
        \hline
    \end{tabular}
    \caption{caption}
    \label{tab:reference}
\end{table}